\documentclass[10pt,a4paper]{amsart}
\usepackage[margin=19mm]{geometry}
\usepackage{amsmath,amssymb,mathtools}
\usepackage[scr=rsfs,cal=euler]{mathalfa}
\usepackage{xfrac}
\usepackage[unicode,hidelinks,bookmarks=false]{hyperref}
\newcommand{\ie}{i.e.\ }
\newcommand{\cf}{cf.\ }
\newcommand{\resp}{respectively\ }
\newcommand{\Subsection}{\subsection}
\providecommand{\nobreakdash}{\nobreak\hbox{-}}
\setlength{\emergencystretch}{2em}
\multlinegap0pt
\usepackage{amssymb}
\usepackage{xcolor}
\usepackage{algorithm}\usepackage{algpseudocode}
\usepackage{cancel}
\usepackage{mathtools}
\newtheorem{proposition}{Proposition}
\def\Ga{\Gamma}
\def\bal#1\eal{\begin{align}#1\end{align}}
\def\bfV{\mathbf{V}}
\def\bfu{\mathbf{u}}
\def\bfv{\mathbf{v}}
\def\bfxi{\boldsymbol{\xi}}
\def\bq{\begin{equation}}
\def\eq{\end{equation}}
\def\md{\mathrm{d}}
\def\nn{\nonumber}
\title{Global existence of weak solutions to incompressible anisotropic Cahn-Hilliard-Navier-Stokes system}
\author{Azeddine Zaidni, Saad Benjelloun, Radouan Boukharfane}
\date{Source: arXiv:2412.05757v4 (CC BY 4.0)}
\begin{document}
\maketitle

\noindent\emph{Source and licence.} Global existence of weak solutions to incompressible anisotropic Cahn-Hilliard-Navier-Stokes system. Version arXiv:2412.05757v4 (CC BY 4.0), distributed by arXiv under the Creative Commons Attribution 4.0 International licence, http://creativecommons.org/licenses/by/4.0/, as recorded in the arXiv licence metadata for this exact version and shown on the abstract page https://arxiv.org/abs/2412.05757v4. Full licence notice: \texttt{licenses/arxiv-2412.05757-CC-BY-4.0.txt}.\par\smallskip

\noindent\textit{Editorial scope.} Counted unit: the article's proposition prop2 (source0.tex lines 805--811). The article's complete original proof of it is delivered with it (source0.tex lines 812--862, one \texttt{proof} environment of 51 lines). No mathematics is written, repaired or re-derived: the statement and the proof are the article's own lines, in the article's own notation, and no retained line is substituted or re-flowed.  Retained uncounted as the source-local context the proof needs: lines 560--580.  Nothing was omitted from any retained span, so the package carries no omission record.  No retained line contains a citation, so the wrapper carries no bibliography and no bibliography entry.  No retained line contains a figure environment, an image-inclusion command or drawing code, so no image is delivered and the package declares no asset dependency.  Editorial formatting only, in addition: an AMS \texttt{amsart} preamble that restates the article's own declarations and macros used by the retained lines, namely the environments \texttt{proposition} on separate counters, together with the standard packages those lines need.  The retained environments print in this document as Proposition 1 in order of appearance, whereas the article prints the same items with the numbering of its own full document.  The complete native source of the article as distributed in the arXiv e-print for this exact version is bundled unchanged as \texttt{sources/169782/source0.tex}.
\begin{proposition}{(Local existence of an approximate solution)}. Given the initial data $(\rho_{0n},\bfu_{0n},\phi_{0n},\mu_{0n})$ constructed as above, there exists a time interval $[0,\widetilde{T}]$ with $\widetilde{T} > 0$ and $(\rho_n^{\varepsilon},\bfu_n^{\varepsilon},\phi_n^{\varepsilon},\mu_n^{\varepsilon})$ such that
\bq
\rho_n^{\varepsilon} \in \mathcal{C}^1(\overline{\Omega}\times[0,\widetilde{T}]),\,\,\bfu_n^{\varepsilon}\in \mathcal{C}^1([0,\widetilde{T}];\bfV_n),\,\, \phi_n^{\varepsilon}\in \mathcal{C}^1([0,\widetilde{T}];H_n), \,\, \mu_n^{\varepsilon}\in \mathcal{C}([0,\widetilde{T}];H_n),
\eq
\bq
\partial_t\rho_n^{\varepsilon} + \bfu_n^{\varepsilon}\cdot\nabla\rho_n^{\varepsilon} = 0 \text{ in } \Omega\times(0,\widetilde{T}),\label{density_T}
\eq
\bq
\begin{aligned}
(\rho_n^{\varepsilon}\partial_t \bfu_n^{\varepsilon},\bfv) + (\rho_n^{\varepsilon}(\bfu_n^{\varepsilon}\cdot\nabla)\bfu_n^{\varepsilon},\bfv) + (\nu(\phi_n^{\varepsilon})\mathbb{D}\bfu_n^{\varepsilon},\nabla\bfv)=& (\rho_n^{\varepsilon}\mu_n^{\varepsilon}\nabla\phi_n^{\varepsilon},\bfv)\\
&-(\rho_n^{\varepsilon}\nabla(F_{\varepsilon}(\phi_n^{\varepsilon})),\bfv),\label{velocity_T}
\end{aligned}
\eq
\bq
(\rho_n^{\varepsilon}\partial_t\phi_n^{\varepsilon},\omega) + (\rho_n^{\varepsilon}\bfu_n^{\varepsilon}\cdot\nabla\phi_n^{\varepsilon},\omega)+ (D(\phi_n^{\varepsilon})\nabla\mu_n^{\varepsilon},\nabla\omega) = 0,\label{concentration_T}
\eq
\bq
(\rho_n^{\varepsilon}\mu_n^{\varepsilon},\omega) =(\Ga\bfxi(\nabla\phi_n^{\varepsilon}),\nabla\omega) + (\rho_n^{\varepsilon} F_{\varepsilon}^{\prime}(\phi_n^{\varepsilon}),\omega),\label{chemical_T}
\eq
for all $\bfv\in \bfV_n$, $\omega\in H_n$ and for all $t\in [0,\widetilde{T}]$.
\end{proposition}

\begin{proposition}\label{prop2}
Given the initial data $(\rho_{0n},\bfu_{0n},\phi_{0n},\mu_{0n})$ and let $(\rho_{n}^{\varepsilon},\bfu_{n}^{\varepsilon},\phi_{n}^{\varepsilon},\mu_{n}^{\varepsilon})$ denote a solution of \eqref{density_T}-\eqref{chemical_T} defined in $[0,T_0)$. There is a constant
$K_0$, depending on $(\rho_{*},\rho^{*},\|\bfu_{0}\|_{L^2(\Omega)},\|\phi_{0}\|_{H^1(\Omega)})$, but independent of $T_0$ , such that 
\bq
\|\bfu_{n}^{\varepsilon}\|_{\mathcal{C}([0,{T}_0];L^2(\Omega))} + \|\phi_{n}^{\varepsilon}\|_{\mathcal{C}([0,{T}_0];H^1(\Omega))}\leq K_0^2.
\eq
\end{proposition}

\begin{proof}\label{proof_prop_2}
Multiplying \eqref{velocity_T} by $\alpha_i^{\varepsilon}$ and summing over $i$ we find
\bq
\begin{aligned}
\frac{ \md}{ \md t}\int_{\Omega} \frac{1}{2}\rho_n^{\varepsilon} |\bfu_n^{\varepsilon}|^2 \md x + \int_{\Omega}\nu({\phi}_n^{\varepsilon})|\mathbb{D}(\bfu_n^{\varepsilon})|^2 \md x=\int_{\Omega}& \rho_n^{\varepsilon} \mu_n^{\varepsilon}\bfu_n^{\varepsilon}\cdot\nabla{\phi}_n^{\varepsilon}\md x\\
&-\int_{\Omega} \rho_n^{\varepsilon} \bfu_n^{\varepsilon}\cdot\nabla F_{\varepsilon}(\phi_n^{\varepsilon})\md x.\label{KIN3}
\end{aligned}
\eq
Now, we multiply
\eqref{concentration_T} by $\gamma_i^{\varepsilon}$ and \eqref{chemical_T} by $\frac{\md }{\md t} \beta_i^{\varepsilon}$, respectively, and summing over $i$, we get
\bq 
\begin{aligned}
\frac{\md }{\md t} \int_{\Omega} \left[\frac{1}{2}\Ga^2(\nabla\phi_n^{\varepsilon}) + \rho_n^{\varepsilon} F_{\varepsilon}(\phi_n^{\varepsilon})\right] \md x &+ \int_{\Omega}\rho_n^{\varepsilon} \bfu_n^{\varepsilon}\cdot\nabla{\phi}_n^{\varepsilon}\mu_n^{\varepsilon}\md x &+ \int_{\Omega}D(\phi_n^{\varepsilon})|\nabla\mu_n^{\varepsilon}|^2 \md x\\
&= \int_{\Omega} \partial_t \rho_n^{\varepsilon} F_{\varepsilon}(\phi_n^{\varepsilon})\md x. \label{Pot_1}
\end{aligned}
\eq
By summing the previous equation \eqref{KIN3} and \eqref{Pot_1}, we obtain 

\bal
\frac{\md}{\md t} \int_{\Omega} \frac{1}{2} \rho_n^{\varepsilon}\left|\bfu_n^{\varepsilon}\right|^2+\frac{1}{2}\Ga^2\left(\nabla \phi_n^{\varepsilon}\right)&+\rho_n^{\varepsilon} F_{\varepsilon}\left(\phi_n^{\varepsilon}\right) \md x\nn\\
&+\int_{\Omega} \nu\left(\phi_n^{\varepsilon}\right)\left|\mathbb{D} \bfu_n^{\varepsilon}\right|^2+D(\phi_n^{\varepsilon})\left|\nabla \mu_n^{\varepsilon}\right|^2 \md x=0 .
\eal

Integrating in time over $[0, t]$ where $0<t<T_0$, we deduce that
\bq
\begin{aligned}
\int_{\Omega}& \frac{1}{2} \rho_n^{\varepsilon}(t)\left|\bfu_n^{\varepsilon}(t)\right|^2 +\frac{1}{2}\Ga^2\left(\nabla \phi_n^{\varepsilon}(t)\right)+\rho_n^{\varepsilon}(t) F_{\varepsilon}\left(\phi_n^{\varepsilon}(t)\right) \md x\\
&+\int_0^t \int_{\Omega} \nu\left(\phi_n^{\varepsilon}(\tau)\right)\left|\mathbb{D} \bfu_n^{\varepsilon}(\tau)\right|^2 \md x \md \tau +\int_0^t \int_{\Omega}D(\phi_n^{\varepsilon}(\tau))\left|\nabla \mu_n^{\varepsilon}(\tau)\right|^2 \md x \md \tau\\
&=\int_{\Omega} \frac{1}{2} \rho_{0 n}\left|\bfu_{0 n}\right|^2+\frac{1}{2}\Ga^2\left(\nabla \phi_{0 n}\right)+\rho_{0 n} F_{\varepsilon}\left(\phi_{0 n}\right) \mathrm{d} x.\label{Enrg_n}
\end{aligned}
\eq

Using the properties of the projector operator $P_n^1$ and $P_n^2$, we find
$$
\begin{aligned}
\int_{\Omega} \frac{1}{2} &\rho_{0 n}\left|\bfu_{0 n}\right|^2  +\frac{1}{2}\Ga^2\left(\nabla \phi_{0 n}\right)+\rho_{0 n} F_{\varepsilon}\left(\phi_{0 n}\right) \mathrm{d} x \\& \leq \frac{\rho^*}{2}\|\bfu_0\|_{L^2(\Omega)}^2+\frac{R}{2}\|\phi_0\|_{H^1(\Omega)}^2+K({\varepsilon}) \rho^*\left(1+\|\phi_0\|_{H^1(\Omega)}^4\right),
\end{aligned}
$$
where the constant $K({\varepsilon}) $ is independent of $n$. Since $F_{\varepsilon}(s) \geq 0$, we obtain
$$
\begin{aligned}
\frac{\rho_*}{2}\|\bfu_n^{\varepsilon}(t)\|_{L^2(\Omega)}^2 +\frac{r}{2}\|\nabla \phi_n^{\varepsilon}(t)\|_{L^2(\Omega)}^2 \leq \frac{\rho^*}{2}\|\bfu_0\|_{L^2(\Omega)}^2&+\frac{R}{2}\|\phi_0\|_{H^1(\Omega)}^2\\
&+K({\varepsilon}) \rho^*\left(1+\|\phi_0\|_{H^1(\Omega)}^4\right).
\end{aligned}
$$

Therefore, there exists a positive constant $K_0=K_0\left({\varepsilon},\rho_*, \rho^*, \bfu_0, \phi_0,r,R\right)$ such that
$$
\sup _{[0 ,T_0]}\|\bfu_n^{\varepsilon}(t)\|_{L^2(\Omega)}+\sup_{[0 ,T_0]}\|\phi_n^{\varepsilon}(t)\|_{H^1(\Omega)} \leq K_0^2.
$$
\end{proof}

\end{document}
